\documentclass[11pt]{article}
\usepackage[margin=1in]{geometry}
\usepackage[T1]{fontenc}
\usepackage{lmodern,amsmath,amssymb,amsthm,booktabs,microtype,xcolor}
\usepackage[colorlinks=true,urlcolor=blue!50!black,linkcolor=blue!50!black]{hyperref}
\newtheorem{theorem}{Theorem}
\newtheorem{lemma}[theorem]{Lemma}
\newtheorem{proposition}[theorem]{Proposition}
\theoremstyle{remark}
\newtheorem*{remark}{Remark}
\hypersetup{pdftitle={Testing a ladder for a rainbow perfect matching: three regimes}}
\title{\LARGE\bfseries Testing a ladder for a rainbow perfect matching:\\ three regimes}
\author{}
\date{Working note \(\,\cdot\,\) November 16, 2026}
\begin{document}
\maketitle

\begin{abstract}\noindent
Fix a list of perfect matchings of the ladder \(L_k=P_k\square K_2\) before
seeing an edge colouring. The list is a \emph{test family for deficit \(d\)}
if, for every colouring with exactly \(|E(L_k)|-d\) colours that has a
rainbow perfect matching, some matching in the list is rainbow. Let
\(h_d(k)\) be the least size of such a list. We show that \(h_d(k)\) is
nondecreasing in \(d\) and passes through three regimes. For \(d\le 2\) it
equals \(d+1\) for every \(k\ge3\), and the optimal families for \(d=2\) are
counted by the Jacobsthal number \(J_{k-2}=(2^{k-2}-(-1)^{k})/3\). For each
fixed \(d\ge3\) it has order \(\log k\). For \(k-1\le d\le 2k-2\), and for no
smaller \(d\), every perfect matching must be tested.
\end{abstract}

\section{Setting}
The ladder \(L_k\) has columns \(0,\dots,k-1\), a rung \(r_i\) in column
\(i\), and for each \emph{gap} \(g\in\{0,\dots,k-2\}\) an upper rail \(u_g\)
and a lower rail \(\ell_g\) joining columns \(g\) and \(g+1\). It has
\(m=3k-2\) edges. A perfect matching is determined by its set \(J\) of gaps:
it contains both rails of each gap in \(J\) and the rungs of the columns
those gaps leave uncovered. The sets \(J\) that occur are exactly the sets
with no two consecutive gaps, so there are \(F_{k+1}\) perfect matchings
(\(F_1=F_2=1\)). Write \(M_J\) for the matching with gap set \(J\). Column
\(c\) is \emph{free} in \(M_J\) if \(r_c\in M_J\), that is, if neither
\(c-1\) nor \(c\) lies in \(J\).

A colouring of deficit \(d\) is an edge colouring with exactly \(m-d\)
colours; equivalently its colour classes \(C\) satisfy \(\sum(|C|-1)=d\). A
matching is rainbow when it meets every class at most once. Since a perfect
matching has \(k\) edges, a rainbow perfect matching can exist only when
\(d\le 2k-2\), and we assume \(k\ge3\) and \(0\le d\le 2k-2\) throughout. We
use one earlier result from this repository, proved in the note
\emph{matching foundations}: the anti-Ramsey value \(A(L_k)=3k-5\) for
\(k\ge3\). In the present language, every colouring of deficit at most two
has a rainbow perfect matching.

\begin{lemma}[Monotonicity]\label{mono}
If \(d<2k-2\) then \(h_d(k)\le h_{d+1}(k)\).
\end{lemma}
\begin{proof}
Let a list fail at deficit \(d\): some colouring has a rainbow perfect
matching \(M\) and no rainbow matching in the list. It has
\(m-d>k\) colours, so some class contains no edge of \(M\). Merge that class
with any other class. The deficit becomes \(d+1\), \(M\) is still rainbow,
and no matching has become rainbow. So the list fails at deficit \(d+1\).
\end{proof}

\section{Small deficit: \texorpdfstring{\(h_d(k)=d+1\)}{h = d+1}}
Clearly \(h_0(k)=1\). For \(d=1\) a single matching \(M\) fails on a
colouring whose only repeated colour lies on two edges of \(M\), while two
matchings with at most one common edge succeed; for instance
\(M_\varnothing\) and \(M_{\{0,2,4,\dots\}}\). So \(h_1(k)=2\).

\begin{lemma}\label{criterion}
Three distinct perfect matchings \(M_1,M_2,M_3\) of \(L_k\) form a test
family for deficit two if and only if
\begin{enumerate}\itemsep0pt
\item[\textup{(a)}] \(|M_i\cap M_j|\le1\) for all \(i\neq j\), and
\item[\textup{(b)}] there are no three distinct edges \(a\in M_1\cap M_2\),
\(b\in M_1\cap M_3\), \(c\in M_2\cap M_3\).
\end{enumerate}
No family of two matchings is a test family for deficit two.
\end{lemma}
\begin{proof}
A colouring of deficit two has either one class \(\{a,b,c\}\) of size three
or two disjoint classes \(P,Q\) of size two. It always has a rainbow perfect
matching, so the family fails exactly when all its members are non-rainbow
for some such colouring.

If (a) fails, say \(P\subseteq M_1\cap M_2\) with \(|P|=2\). If
\(P\subseteq M_3\) take any pair \(Q\) disjoint from \(P\); otherwise
\(M_3\setminus P\) has at least \(k-1\ge2\) edges and we take \(Q\) inside
it. Every \(M_i\) contains \(P\) or \(Q\). If (b) fails, the class
\(\{a,b,c\}\) meets each \(M_i\) twice.

Conversely assume (a) and (b). With classes \(P,Q\), two of the three
matchings would contain the same pair, against (a). With a class
\(\{a,b,c\}\), each \(M_i\) contains at least two of the three edges; by (a)
no two matchings contain the same two, so up to renaming
\(M_1\supseteq\{a,b\}\), \(M_2\supseteq\{a,c\}\), \(M_3\supseteq\{b,c\}\),
against (b).

For two matchings \(M_1\ne M_2\), choose \(e\in M_1\setminus M_2\) and
another \(f\in M_1\), put \(P=\{e,f\}\), and take \(Q\) inside
\(M_2\setminus P\), which has at least \(k-1\) edges. Both matchings fail.
\end{proof}

\begin{theorem}\label{three}
For every \(k\ge3\), \(h_2(k)=3\). The test families of size three for
deficit two are exactly the triples \(\{M_A,M_B,M_C\}\) in which \(A,B,C\)
partition the set of all gaps and the first and last gap lie in different
parts. Their number is the Jacobsthal number
\[
  J_{k-2}=\frac{2^{k-2}-(-1)^{k-2}}{3}\qquad(1,1,3,5,11,21,43,\dots\ \text{for }k=3,4,5,\dots).
\]
\end{theorem}
\begin{proof}
Let \(A,B,C\) be the gap sets of a test family of size three. Two matchings
that share a gap share its two rails, so by (a) the sets are pairwise
disjoint, and the only possible common edges are rungs of columns free in
both.

Column \(0\) is covered only by gap \(0\), which lies in at most one set, so
it is free in at least two of the matchings; likewise column \(k-1\). If
column \(0\) were free in all three, then \(r_0\) would lie in every
pairwise intersection, and a pair in which column \(k-1\) is free would
share two rungs. So gap \(0\) lies in exactly one set, and similarly gap
\(k-2\). They lie in different sets: otherwise the other two matchings share
\(r_0\) and \(r_{k-1}\). Name the sets so that \(0\in A\) and \(k-2\in B\).

Now \(r_0\in M_B\cap M_C\) and \(r_{k-1}\in M_A\cap M_C\). An edge of
\(M_A\cap M_B\) would differ from both, since \(r_0\notin M_A\) and
\(r_{k-1}\notin M_B\), and (b) would fail. Hence \(M_A\cap M_B=\varnothing\),
and by (a) the other two intersections are exactly \(\{r_0\}\) and
\(\{r_{k-1}\}\). In terms of columns: no interior column is free in two of
the matchings. So each of the \(k-2\) interior columns is covered in at
least two of the matchings, and columns \(0\) and \(k-1\) are covered once
each. Counting covered (column, matching) incidences,
\[
  2\bigl(|A|+|B|+|C|\bigr)\ \ge\ 2(k-2)+2 .
\]
As the sets are disjoint subsets of the \(k-1\) gaps, equality holds and
they partition the gaps.

Conversely, let \(A,B,C\) partition the gaps, each without consecutive gaps,
with \(0\in A\) and \(k-2\in B\). The two gaps beside an interior column lie
in different parts, so that column is free in exactly one matching. The
pairwise intersections are therefore \(\varnothing\), \(\{r_0\}\) and
\(\{r_{k-1}\}\), and (a), (b) hold.

Such partitions are the proper colourings of the path on the \(k-1\) gaps
with three colours and differently coloured ends, taken up to renaming the
colours. Renaming acts freely, and the number of proper colourings with
prescribed distinct end colours is the number of walks of length \(k-2\)
between two given vertices of a triangle, which is \((2^{k-2}-(-1)^{k-2})/3\).
The lower bound \(h_2(k)\ge3\) is the last part of Lemma~\ref{criterion}.
\end{proof}

\begin{remark}
In an optimal family every rail and every interior rung is tested exactly
once, and each end rung exactly twice. For \(d=1\) the optimal pairs are the
pairs with at most one common edge; there are \(3,5,13,25,57,117,253\) of
them for \(k=3,\dots,9\) (computed).
\end{remark}

\section{Fixed deficit at least three: order \texorpdfstring{\(\log k\)}{log k}}
\begin{theorem}\label{log}
Let \(k\ge5\) and \(3\le d\le 2k-2\). Then
\[
  \Bigl\lceil\log_3\bigl\lfloor (k-1)/2\bigr\rfloor\Bigr\rceil\ \le\ h_d(k)
  \ \le\ \Bigl\lfloor 8^{d}\ln\tbinom{3k-2}{d}\Bigr\rfloor+1 .
\]
In particular \(h_d(k)=\Theta_d(\log k)\) for each fixed \(d\ge3\), in
contrast with \(h_2(k)=3\).
\end{theorem}
\begin{proof}
\emph{Lower bound.} By Lemma~\ref{mono} it is enough to treat \(d=3\). Use
the odd columns from \(1\) to \(k-2\). In a perfect matching such a column
is in one of three states: free, covered by the gap on its left, or covered
by the gap on its right. A list of \(t\) matchings gives each column a word
of length \(t\). If odd interior columns \(i<j\) have the same word, give
each of
\[
 \{r_i,r_j\},\qquad\{u_{i-1},u_{j-1}\},\qquad\{u_i,u_j\}
\]
its own repeated colour and colour all other edges distinctly. Since
\(j\ge i+2\) the three pairs are disjoint. A listed matching puts columns
\(i\) and \(j\) in the same state and so contains one of the pairs, while
\(M_{\{i-1\}}\) contains only \(u_{i-1}\) and \(r_j\) among the six edges
and is rainbow. Hence the \(\lfloor(k-1)/2\rfloor\) words are distinct and
\(3^t\ge\lfloor(k-1)/2\rfloor\).

\emph{Upper bound.} Given a colouring of deficit \(d\) with a rainbow
perfect matching \(M\), keep one edge of each colour, taking the edge of
\(M\) when there is one. The set \(F\) of discarded edges has size \(d\),
\(M\) avoids it, and every perfect matching avoiding \(F\) is rainbow. So it
suffices that for every \(d\)-set \(F\) avoided by some perfect matching, a
listed matching avoids \(F\).

Choose independent fair bits on the gaps. Select the gaps with bit one whose
neighbouring bits are both zero; these are pairwise nonconsecutive. Fix
\(F\) and a matching \(M_J\) avoiding it. For a rung \(r_c\in F\), column
\(c\) is covered by a gap of \(J\): ask for bit one there and zero on its
neighbours. For a rail in \(F\), its gap is not in \(J\): ask for bit zero
there. These at most \(3d\) requirements agree with the indicator of \(J\),
so they hold together with probability at least \(8^{-d}\), and then the
random matching avoids \(F\). With \(t\) independent matchings the expected
number of missed sets \(F\) is at most
\(\binom{3k-2}{d}(1-8^{-d})^t<\binom{3k-2}{d}e^{-t/8^d}\), which is less
than one for the stated \(t\).
\end{proof}

\section{Large deficit: every matching is needed}
\begin{theorem}\label{sat}
For \(k\ge3\), \(h_d(k)=F_{k+1}\) if and only if \(k-1\le d\le 2k-2\).
\end{theorem}
\begin{proof}
Let \(d=k-1\) and fix a gap set \(J\). For each gap \(g\in J\) make
\(\{r_g,r_{g+1}\}\) a colour class, for each gap \(g\notin J\) make
\(\{u_g,\ell_g\}\) a class, and colour the other edges distinctly. The
\(k-1\) classes are disjoint, so the deficit is \(k-1\), and \(M_J\) meets
none of them twice. Let \(J'\ne J\). If \(J'\) has a gap outside \(J\), then
\(M_{J'}\) contains both rails of it. Otherwise \(J'\subsetneq J\); take
\(g\in J\setminus J'\). Gaps \(g-1\) and \(g+1\) are not in \(J\), hence not
in \(J'\), so columns \(g\) and \(g+1\) are free in \(M_{J'}\) and it
contains \(r_g,r_{g+1}\). Thus \(M_J\) is the only rainbow perfect matching
and must be in every test family. This holds for all \(J\), so
\(h_{k-1}(k)=F_{k+1}\), and Lemma~\ref{mono} extends it to
\(k-1\le d\le2k-2\).

Let \(d\le k-2\). We claim the \(F_{k+1}-1\) matchings other than
\(M_\varnothing\) form a test family. If \(M_\varnothing\) is rainbow, keep
one edge of each colour including all rungs. The \(d\) discarded edges are
rails, and there are \(k-1>d\) gaps, so some gap \(g\) keeps both rails.
Then \(M_{\{g\}}\) consists of kept edges and is rainbow. If \(M_\varnothing\) is not rainbow, every rainbow perfect matching is already in the family.
\end{proof}

\section{Exact values for short ladders}
An exhaustive search over all colourings of the given deficit, followed by
an exact minimum hitting-set computation, gives the following.
\begin{center}\small
\begin{tabular}{@{}rrrrrrrr@{}}
\toprule
\(k\) & \(F_{k+1}\) & \(d=0\) & \(1\) & \(2\) & \(3\) & \(4\) & \(5\)\\
\midrule
3 & 3  & 1 & 2 & 3 & 3 & 3 & --\\
4 & 5  & 1 & 2 & 3 & 5 & 5 & 5\\
5 & 8  & 1 & 2 & 3 & 7 & 8 & 8\\
6 & 13 & 1 & 2 & 3 & 7 & 12 & 13\\
7 & 21 & 1 & 2 & 3 & 7 & 15 &\\
8 & 34 & 1 & 2 & 3 & 8 & &\\
9 & 55 & 1 & 2 & 3 & 8 & &\\
\bottomrule
\end{tabular}
\end{center}
Blank entries were not computed. The table agrees with Theorems~\ref{three} and~\ref{sat} and with
monotonicity. The number of optimal families at \(d=2\) found by the search
is \(1,1,3,5,11,21,43\) for \(k=3,\dots,9\).

\section{Questions}
The bounds in Theorem~\ref{log} are far apart: the table suggests that
\(h_3(k)\) is larger than the lower bound and grows slowly. Is
\(h_3(k)/\log_2 k\) convergent, and what is the dependence of \(h_d(k)\) on
\(d\) in the range \(3\le d\le k-2\)? For \(k=5,6\) the table gives
\(h_{k-2}(k)=F_{k+1}-1\); whether this persists is open.
Finally, the definitions make sense for any graph with a perfect matching;
the cycle ladder (prism) and the grid \(P_k\square P_3\) are natural next
cases.

\bigskip\small\noindent
\textbf{Status, method and related work.}
This is a working note that has not been externally refereed. The statements
were found and proved with substantial assistance from AI systems and were
checked by the computations described above; the accompanying script
\texttt{experiments/ladder\_thresholds.py} re-derives the counts in
Theorem~\ref{three} by brute force over all deficit-two colourings for
\(k\le7\) and checks the colourings used in Theorems~\ref{log}
and~\ref{sat} for \(k\le12\). The upper bound in Theorem~\ref{log} is a
standard random-restriction argument in the spirit of the universal sets of
Naor, Schulman and Srinivasan (\emph{Splitters and near-optimal
derandomization}, FOCS 1995). Anti-Ramsey numbers of matchings have been
determined for complete graphs (Haas and Young; \"Ozkahya and Young), for
3-regular bipartite hosts (Jin 2017; a 2024 paper in \emph{Discrete
Mathematics} obtains \(3m-5\) for 3-edge-connected cubic bipartite graphs on
\(2m\) vertices, the same value as \(A(L_k)\) although ladders are not
cubic), and for outerplanar families. A search on November 16, 2026 did not
find the test-family parameter \(h_d\) or the results above in the
literature, but that search was limited to open web sources and does not
establish novelty.
\end{document}
